% RESPONSAVEL: gerencia (Guilherme e colega)
\section{Introdu\c{c}\~ao ao Problema}

O Mundo dos Blocos \'e um dom\'inio cl\'assico de planejamento em Intelig\^encia Artificial.
Dado um estado inicial e um estado objetivo, o problema \'e encontrar uma sequ\^encia de a\c{c}\~oes que leve de um ao outro.
Na vers\~ao cl\'assica, todos os blocos t\^em o mesmo tamanho e cada bloco ocupa exatamente uma posi\c{c}\~ao.
Neste trabalho estendemos o dom\'inio para blocos de \emph{comprimentos diferentes}, dispostos sobre uma mesa com pontos discretos e empilh\'aveis em n\'iveis.

\subsection{O dom\'inio}

Temos quatro blocos $a$, $b$, $c$ e $d$, com comprimentos $\ell(a)=\ell(b)=1$, $\ell(c)=2$ e $\ell(d)=3$.
A mesa \'e um segmento com os pontos inteiros de $0$ a $6$, o que d\'a seis \emph{slots} unit\'arios, $s_i=[i,i+1]$ para $i\in\{0,\dots,5\}$.
Um bloco come\c{c}a em um ponto $p$ e cobre $\ell(b)$ slots consecutivos. Os blocos podem ficar sobre a mesa (n\'ivel $0$) ou sobre outros blocos (n\'iveis acima).

Essa extens\~ao traz quatro caracter\'isticas que o caso simples n\~ao tem:
\begin{itemize}
\item \textbf{Tamanhos diferentes:} um bloco pode cobrir mais de um slot e se apoiar em mais de um bloco.
\item \textbf{Espa\c{c}os vazios:} o bloco do n\'ivel de cima pode ficar em ponte sobre dois blocos de baixo, deixando um slot vazio entre eles.
\item \textbf{Movimentos laterais:} al\'em de subir e descer, um bloco pode deslizar na horizontal, ent\~ao a posi\c{c}\~ao em que ele \'e colocado importa.
\item \textbf{Equil\'ibrio:} um bloco maior sobre um menor s\'o \'e aceito se tiver apoio suficiente por baixo. A regra adotada exige apoio em pelo menos $\lceil \ell(b)/2\rceil$ dos slots que o bloco cobre.
\end{itemize}

\subsection{Objetivo do trabalho}

O objetivo \'e propor uma representa\c{c}\~ao do conhecimento para esse dom\'inio e us\'a-la em racioc\'inio automatizado para gerar planos.
O trabalho segue as etapas abaixo, na ordem em que o enunciado pede.
\begin{enumerate}
\item Representar o dom\'inio em L\'ogica de Primeira Ordem (LPO), com os predicados, as leis do dom\'inio e a a\c{c}\~ao $\Move$.
\item Indicar, para cada elemento da representa\c{c}\~ao, o que a a\c{c}\~ao adiciona ou remove (adds e deletes).
\item Usar a representa\c{c}\~ao para obter, manualmente, os planos de cada cen\'ario.
\item Considerar ordem parcial entre as metas, na forma $A \xrightarrow{P} B$.
\item Somente depois, codificar a solu\c{c}\~ao de LPO para l\'ogica proposicional (CNF), gerar os planos com um SAT solver e compar\'a-los com os planos manuais.
\end{enumerate}

\subsection{As tr\^es situa\c{c}\~oes}

O enunciado define tr\^es situa\c{c}\~oes sobre o mesmo conjunto de blocos.
Na Situa\c{c}\~ao 1, parte-se do estado inicial $S_0$ e o plano deve chegar a um dos estados objetivo $S_{f1}$, $S_{f2}$, $S_{f3}$ ou $S_{f4}$.
Na Situa\c{c}\~ao 2, o plano vai de $S_0$ at\'e $S_5$.
Na Situa\c{c}\~ao 3, o plano vai de $S_0$ at\'e $S_7$.
Em todas elas a \'unica a\c{c}\~ao dispon\'ivel \'e $\Move(b,y,p)$, que move o bloco $b$ para cima de $y$ (outro bloco ou a mesa $T$) come\c{c}ando no ponto $p$.

\subsection{Organiza\c{c}\~ao do documento}

A Se\c{c}\~ao 2 apresenta a descri\c{c}\~ao formal do dom\'inio (itens 1, 2 e 4) e a execu\c{c}\~ao manual das tr\^es situa\c{c}\~oes (item 3).
As Se\c{c}\~oes 3 a 7 tratam da codifica\c{c}\~ao em CNF e do uso do SAT solver (item 5): as vari\'aveis e cl\'ausulas, os tr\^es cen\'arios codificados, o mapeamento entre a descri\c{c}\~ao formal e o c\'odigo, a execu\c{c}\~ao passo a passo e a interpreta\c{c}\~ao da sa\'ida do solver.
